% from ega4-18-fr.tex
\oldpage[IV-4]{161}

  \item[(iii)] Si l'ensemble des composantes irréductibles de $Y$ est localement fini, il en est de même de l'ensemble des composantes irréductibles de $X$, et en particulier $X$ est localement connexe.
\end{enumerate}
\end{lemma}

L'assertion (i) résulte aussitôt de (2.3.4) et ($0_{\mathrm I}$, 2.1.6) et du fait que les fibres $f^{-1}(y)$ sont des ensembles discrets. Il est clair que (ii) est un cas particulier de (i). Enfin, pour prouver (iii), on peut, en vertu de (i), se borner au cas où $f$ est quasi-fini et où $Y$ est irréductible ($0_{\mathrm I}$, 2.1.6) et la conclusion résulte alors de (ii) et du fait que les fibres de $f$ sont des ensembles finis (II, 6.2.2).

\begin{proposition}[18.10.7]
\label{IV.18.10.7}
Soient $Y$ un préschéma géométriquement unibranche et irréductible (resp. intègre), $f:X\to Y$ un morphisme étale. Alors $X$ est isomorphe à une somme de préschémas irréductibles (resp. intègres) $X_\lambda$, où $\lambda$ parcourt la fibre $f^{-1}(\eta)$ du point générique $\eta$ de $Y$. Les $X_\lambda$ sont géométriquement unibranches; si $Y$ est normal, les $X_\lambda$ sont normaux.
\end{proposition}

En effet, $X$ est géométriquement unibranche (17.5.7), donc tout point de $X$ n'appartient qu'à une seule composante irréductible; en outre, en vertu de (18.10.6) et (17.6.1) l'ensemble des composantes irréductibles de $X$ est localement fini, donc ($0_{\mathrm I}$, 2.1.6) ce sont les composantes connexes $X_\lambda$ de $X$ ($\lambda\in f^{-1}(\eta)$), et $X$ est somme des $X_\lambda$; en vertu de (17.5.7), si $Y$ est intègre (resp. normal), les $X_\lambda$ sont intègres (resp. normaux).

\begin{proposition}[18.10.8]
\label{IV.18.10.8}
Soient $Y$ un préschéma normal et intègre, de point générique $\eta$, $K=k(\eta)$ son corps de fonctions rationnelles, $f:X\to Y$ un morphisme étale. On suppose que $f^{-1}(\eta)$ est finie et non vide, de sorte que le $K$-schéma $f^{-1}(\eta)$ est le spectre d'une $K$-algèbre finie séparable $L$, composée directe d'un nombre fini de corps $K_i$ ($1\leq i\leq n$), extensions finies séparables de $K$ (17.6.2). Soit $Y'$ la fermeture intégrale de $Y$ dans $L$, isomorphe à la somme des fermetures intégrales $Y_i$ de $Y$ dans les $K_i$ (II, 6.3.6). Alors le morphisme $f$ se factorise d'une seule manière en $f:X\xrightarrow{f'}Y'\xrightarrow{g}Y$, tel que la restriction $f'_\eta:f^{-1}(\eta)\to g^{-1}(\eta)$ soit l'isomorphisme canonique. En outre $f'$ est un isomorphisme local; pour que $f'$ soit une immersion ouverte, il faut et il suffit que $f$ soit séparé.
\end{proposition}

En vertu de (18.10.7), on peut se borner au cas où $X$ est intègre et normal, $L$ son corps de fonctions rationnelles, $f$ étant dominant. L'existence et l'unicité de la factorisation considérée de $f$ résultent alors de (II, 6.3.9). Comme $f'$ est birationnel et localement quasi-fini, les dernières assertions résultent de (8.12.10).

\begin{corollary}[18.10.9]
\label{IV.18.10.9}
Sous les hypothèses de (18.10.8), pour que $f$ soit un morphisme fini (autrement dit, pour que $X$ soit un revêtement étale de $Y$), il faut et il suffit que $X$ soit isomorphe à la fermeture intégrale $Y'$ de $Y$ dans $L$.
\end{corollary}

Si $f$ est fini, l'injection canonique $j:X\to Y'$ (qui est une immersion ouverte) est aussi un morphisme fini (II, 6.1.5, (v)), donc un morphisme fermé (II, 6.1.10), et comme $j(X)$ est dense dans $Y'$ par (18.10.8), on a $j(X)=Y'$. Inversement, comme $Y$ est normal et $L$ composé direct d'extensions finies séparables de $K$, $Y'$ est fini sur $Y$ (Bourbaki, \emph{Alg. comm.}, chap.~V, \textsection~1, n\textsuperscript{o}~6, cor.~1 de la prop.~18), d'où le corollaire.

\begin{env}[18.10.10]
\label{IV.18.10.10}
Soient $Y$ un préschéma normal et intègre, $K=R(Y)$ son corps de fonctions rationnelles. Nous dirons qu'une $K$-algèbre de rang fini $L$ est non ramifiée sur $Y$ si: $1^\circ$ $L$ est une $K$-algèbre séparable, donc composée directe d'extensions finies séparables $K_i$ de $K$ ($1\leq i\leq n$); $2^\circ$ la fermeture intégrale $Y'$ de $Y$ dans $L$ (somme des préschémas fermetures intégrales de $Y$ dans les $K_i$) est non ramifiée sur $Y$ (ce qui, par (18.10.3),

